\section{The global Selmer estimate for a specified cut}

% BEGIN S84G1
\paragraph{S84G1. Integral determinant and coboundaries on an arbitrary support.}
Use the data in $\mathrm{StepsNoCut}.\mathrm{GlobalStatement}$.
Write $X$ for the original points, $U$ for the tested subset,
and $\pi\in U\cap Z$. Put $O=O_i$, $K=\operatorname{Frac}O$,
$p=\ell_i$, $e=\ell_O(O/pO)$. Use the given good congruence
objects (in particular their actual integral algebra and
root cover) and put $\delta=\mathrm{del51}(\pi')$ at the
specified lift. Write $\ell$ for $O$-length in this section.
In the native normalization $\ell(O/\delta O)=e v_\delta$.
For the rank block and each character block write $C_b$
for the span of the supported equivariant cocycles of G79E7
with values in $(K/O)^{d_b}$.

First,
\[
   \ell(O/J_\pi)\le \ell(C_\rho)+\sum_{0\le j<t}\ell(C_j),
   \qquad J_\pi=\{h(\pi):h\in T_i,\ h|_{U\setminus Z}=0\}.
\]
Here is the support modification to G79E11--G79E12.
In the span of all $y_g=(\chi_1(g)^{-1}r(g)-1)E_1$ over
the matrix algebra $A$, identify two vectors if their
difference vanishes at every point above $U$. This quotient
is the image of restriction, a finite $O$-module in the
sense of finite generation. Choose finitely many $y_{g_q}$
generating it over $A$. The relations are now precisely the
rows vanishing on this restriction. As in G79E12 each solution
of the block row equations defines an equivariant map on
this image; injectivity of $K/O$ over the DVR extends it to
the ambient matrix module on the restricted points, hence
to a supported map on the original family. The Smith determinant
argument on evaluations at the $g_q$ gives the same sum of
block lengths, provided its determinant coefficients take
values in $J_\pi$ at $\pi'$.

Indeed those coefficients belong to the given root algebra
$T'$ by the single-block determinant condition and successive
formal block elimination in G79E11. Above $x\in U\setminus Z$
the absence of a $\chi_1$-line follows from $\mathrm{admG6}$:
a nonzero fixed vector in the twist is such a line and would
put $x$ in $Z$. As $E_1(x')$ has rank one, some $y_g(x')$
does not kill some vector $v$. Thus the column
$(y_{g_q}(x')v)_q$ is nonzero, and every relation annihilates
this column. The determinant polynomial vanishes on the
whole fiber. Expand a coefficient in the given free $T_i$-basis
of $T'$. For this fiber multiply the expansion by a Hecke
function supported at $x$ and nonzero there (multiply the
pairwise separating factors). The product vanishes everywhere,
so freeness and the domain property of $O$ force each original
coordinate to vanish at $x$. Thus every coordinate vanishes
on $U\setminus Z$ and has its value at $\pi$ in $J_\pi$.
Evaluation at $\pi'$ and summation prove the required assertion.
This works for every $U$ in the original quantifier.

The total loss in passing from these cocycles to cohomology
is at most $2(r+t-1)e v_\delta$. Indeed the chosen projector
has square $\delta E_1$ at $\pi'$, acts by $\delta$ on the
first line and by zero on the other blocks. Its rank there
is one, so $\delta\ne0$ by the distinguished block formula.
The calculation with its integral group-algebra expansion in
G79E7 and G81S3 applies before division: a coboundary
$g\mapsto(M_b(g)-1)v$ in $C_b$ forces $\delta^2 v=0$ for
every other block. This costs at most $2d_b\ell(O/\delta O)$
on that block. The first line has trivial ratio character
and zero coboundaries. The result holds when $\delta$ is a
nonunit, including at the two type primes.
% END S84G1

% BEGIN S84G2
\paragraph{S84G2. The original global inequality, including its local defects.}
Assume global Artin reciprocity. With the hypotheses and the
subset of S84G1, suppose that for all $w\mid p$ the original
$\mathrm{HfChar}$ and $\mathrm{LocBoundP}$ tests hold with
the stated $e_\rho,e_\chi$. Retain the source finiteness
conditions on the rank Selmer and on the indicated cyclotomic
Tate--Shafarevich lengths. Then
\[
\begin{split}
 \mathrm{vccCut}(i,U,\pi)\ \le{}&
  \#\{w\mid p\}(e_\rho+t e_\chi)+s_\pi
  +\sum_{0<j<t}\mathrm{shaLen}_{p}(j)\\
  &+v_p(h_E)+4(r+t)v_\delta
                        +t\sum_{q\in Q_i}v_\zeta(q).
\end{split}
\]
This is the native $\mathrm{GlobalStatement}$ with its
original arbitrary support and coefficient normalization.

For each $b$ and $w\mid p$, restrict $C_b$ to the decomposition
group and quotient its local span by its intersection with
the finite condition. Its image has length at most
$e e_b$, where $e_b=e_\rho$ for the rank block and $e_\chi$
otherwise. To verify this point, take the simultaneous columns
of $E_1$ over the points of the cover above $U$ and the
module spanned by their $y_g$-columns. Reassembly followed
by the supported equivariant map is well-defined, because it
kills vectors vanishing on this support. The resulting
finite list, vector, stable lattice submodule and map into
the block are a generator in the native local cut span,
by the calculation of G79E13. The coefficient isomorphism
of the local bridge identifies $O$ with the tested integers
and preserves $e$. The quotient of a span by its intersection
with the finite condition is its image modulo that condition.
Taking the kernel of the product of all these restriction
maps therefore costs at most $\#\{w\mid p\}e e_b$.

On the rank block the remaining classes belong to the
specified Selmer group. Outside $p,S_i,Q_i$ the inertia of
$r_x$ is trivial by $\mathrm{admG4}$, and that of $\chi_1$ is
trivial there by the decomposition at $\pi$. Hence $y_g$
vanishes on inertia there. The defining Selmer conditions
allow the remaining $S_i,Q_i$. The length of its image in
classes is at most $e s_\pi$.

On each character block the same unramified assertion holds.
At $S_i$, $\mathrm{admG5}$ and the tame-generator conditions
give either the integral polynomial relation, with its
evaluation at $\chi_1(\gamma_0)$ a unit, or the quadratic
inertia relation and $\zeta-1$ a unit. In the polynomial
case use that relation on $y_g$ and evaluate by the character
block: all the ratio characters are unramified here because
this place is away from $p$. Thus the unit annihilates the
cocycle value. In the quadratic case evaluation on the
tame generator puts the inertia image into $(K/O)[\zeta-1]$,
as in G79E8. Evaluation there detects the whole restriction:
it is additive on inertia, continuous with finite $p$-power
image, and kills the pro-residue-characteristic wild group.
Thus both alternatives give zero inertia restriction at $S_i$.
At $q\in Q_i$ the quadratic relation holds with
$\zeta_i(q)\ne1$. Taking the kernel of all these inertia
restrictions loses at most $e\sum_{q\in Q_i}v_\zeta(q)$
on this character block.

For $j=0$ the first-character clause of $\mathrm{HfChar}$
makes the remaining cocycles unramified also at $p$.
They factor through the ideal class group, with length at
most $e v_p(h_E)$ by G79E6. For $0<j<t$ the ratio is
$\chi_{\rm cyc}^{-j}$. The finite-character clause gives
local coboundaries above $p$. The classes unramified away
from $p$ and locally trivial above it have length at most
$e\,\mathrm{shaLen}_p(j)$ by the coefficient-coordinate
and inflation argument in G79E6. Its kernel consists of
global coboundaries, already bounded with the rank coboundaries
in S84G1. Summing all the block inequalities, applying the
determinant inequality there, and dividing by $e>0$ proves
the display: $2(r+t-1)v_\delta\le4(r+t)v_\delta$, since
$v_\delta\ge0$. The evaluation-ideal identification uses
only that the off-locus is $U\setminus Z$, as in P84D1.
% END S84G2

% BEGIN S84G3
\paragraph{S84G3. Application at the completely split type prime.}
On a good cut take an actual tested type index $p=g_k$, a
good point $\pi$, a tested good integral family $r$, and the
tested rank family $\rho$ with its semisimple decomposition
at $\pi$. Use the Galois hypotheses of C84F1--C84F6.
The point $\pi\in Z_{\rm good}$ witnesses
$\mathrm{HG}$ at this cut index; also $\pi\in Z$. The fourth
coefficient duty of C84F6 applies to exactly this $i,k,r,\rho,\pi$
and this decomposition. Its result
$\mathrm{WildIdempotentChoice}$ consumes $\pi\in Z$ again
and gives a record $S$ satisfying
$\mathrm{GoodCongruenceObjects}(D_{c}.\mathrm{nd},r,\rho,\pi,S)$
and the wild inertia condition for this same $S.g_0$.
In particular it furnishes the precise integral object in
S84G1 on the cut spectrum before the support is chosen for
its determinant calculation. The replacement of the model
family by this particular tested family is justified by the
integral conjugacy argument in C84F6.

For each $\pi$ use the support $U=S_\pi$ of P84D4, with
$\pi\in U$ and the asserted single local type tuple on all
of $U$. Assume the local hypotheses of T84L1--T84L6 for
this support, including the direct potentially crystalline
family there, the stated host, size and jump comparisons,
and the quantitative lifting input $\mathrm{admT}(C_T)$.
Also assume global Artin reciprocity and the native Selmer
and cyclotomic finiteness tests, as in S84G2. For $n=r+t$,
$a=\eta+t+1$, set $N=n+1$ and $u=\lfloor N^2/4\rfloor$.
The actual $e_{\max}=C_T(Nu+a(N+u))$ works for all
these places. T84L6 and the local bounds give
$e_\rho=2r e_{\max}$ and $e_\chi=2e_{\max}$. There are
$[E:\mathbb Q]=2f$ places above the completely split type
prime. Applying S84G2 on $S_\pi$ gives
\[
     c_\pi^*\le 4fn C_T\bigl(Nu+a(N+u)\bigr)+A_\pi
\]
with the full $A_\pi$ of P84D3. Here the arithmetic constraints
$2\le N\le n+1$, $1\le u$ and $4u\le N^2$ follow from
$n\ge3$. This holds simultaneously on the separately chosen
supports under the stated local tests. P84D5 absorbs
$v_p(N_G)$ into its $2fr^2$ reserve and sums the
multiplicities, proving the native paired estimate with
$C=C_T$. If the same direct potentiality and lifting tests
are given on the full cut spectrum, take that whole spectrum
as $U$ in S84G2 to obtain the full-spectrum $c_\pi$ bound
and use P84D3 directly.
% END S84G3
